\clearpage
\SourceLeaf{242}\pdfbookmark[1]{Sectio tertia.}{liber-secundus-sectio-tertia}
\begin{center}
{\large SECTIO TERTIA.}\par\medskip
\emph{Determinatio orbitae observationibus quotcunque quam proxime satisfacientis.}\par\bigskip
\hypertarget{native-242-article172}{}\pdfbookmark[2]{172.}{article-172}172.
\end{center}
Si observationes astronomicae ceterique numeri, quibus orbitarum computus
innititur, absoluta praecisione gauderent, elementa quoque, sive tribus
observationibus sive quatuor superstructa fuerint, absolute exacta statim
prodirent (quatenus quidem motus secundum leges \textsc{Kepler}i exacte fieri
supponitur), adeoque accitis aliis aliisque observationibus confirmari tantum
possent, haud corrigi. Verum enim vero quum omnes mensurationes atque
observationes nostrae nihil sint nisi approximationes ad veritatem, idemque de
omnibus calculis illis innitentibus valere debeat, scopum summum omnium
computorum circa phaenomena concreta institutorum in eo ponere oportebit, ut
ad veritatem quam proxime fieri potest accedamus. Hoc autem aliter fieri
nequit, nisi per idoneam combinationem observationum \emph{plurium}, quam quot
ad determinationem quantitatum incognitarum absolute requiruntur. Hoc negotium
tunc demum suscipere licebit, quando orbitae cognitio approximata iam
innotuit, quae dein ita rectificanda est, ut omnibus observationibus
\emph{quam exactissime} satisfaciat. Etiamsi haec expressio aliquid vagi
implicare videatur, tamen infra principia tradentur, secundum quae problema
solutioni legitimae ac methodicae subiicietur.

Praecisionem summam ambire tunc tantummodo operae pretium esse potest, quando
orbitae determinandae postrema quasi manus apponenda est. Contra quamdiu spes
affulget, mox novas observationes novis correctionibus \SourceLeaf{243}occasionem daturas esse, prout res fert plus minusve ab extrema praecisione
remittere conveniet, si tali modo operationum prolixitatem notabiliter
sublevare licet. Nos utrique casui consulere studebimus.
